\section{Introduction}
\label{sec:intro}

% Framing: measurement -> paradigm test (suppl. §A) -> objective.
% TODO(human): voice pass.
% 3DV pass 2026-07-13: appendix refs -> supplementary; light trim only.
% Advisor round 3 (2026-07-16): allocation-vs-objective framing made
% explicit; research hypotheses RH1-3 named (RH labels avoid colliding
% with the homology classes H0/H1/H2); numbers + % src unchanged.

Differentiable rasterization reconstructs a triangle soup from images by
following photometric gradients~\citep{tojo2026diffsoup}, periodically
resampling triangles under a fixed budget. Photometric losses are
topologically blind: reconstructions that differ in genus, connectivity, or
enclosed-void structure can be indistinguishable to Chamfer and Hausdorff
while the bottleneck distance between persistence diagrams separates them
by an order of magnitude (suppl.\ Table~S1).
% src: figures/summary.md (Phase 1); h2_unified + crossover reports (suppl. §A)

A decade of
connectivity-free reconstruction --- point splats, Gaussians, triangle
soups --- has treated topology errors as symptoms of geometric
under-provision: represent the surface better, or allocate the primitive
budget more wisely, and topology should follow. The allocation half is
directly testable, and we tested it first --- biasing where the
resampler spends triangles in a controlled prior study (suppl.\ \S A).
It fails: a topological prior helps enclosed voids, but any equally-wide
\emph{non-topological} prior reproduces most of the gain, and no prior
shape repairs loops --- concentrating manufactures phantom handles,
spreading never beats a random control.
% src: PHASE2B_CROSSOVER.md; h2_unified/report/results.json

The null result relocates the root cause: not \emph{where} triangles are
spent, but \emph{what the objective measures} --- a photometric loss is
a local, pixel-space signal; a loop or an enclosed void is a global,
categorical property of the surface. This paper therefore moves topology
into the \emph{objective}: a differentiable persistence term in the
training loss, so gradients act on the birth/death values of the
reconstruction's topological features directly, rather than on where
triangles respawn.

Three research hypotheses structure the study, each tested by a
condition designed in advance. \textbf{RH1 (diagnosis)}: topology errors
originate in the training objective, not in primitive allocation ---
tested by the allocation study above. \textbf{RH2 (specificity)}: a
persistence term improves topology in ways no norm-matched
non-topological regularizer through the identical channel can reproduce
--- the control C2 (\S\ref{sec:conditions}). \textbf{RH3
(complementarity)}: objective and allocation are complementary channels,
not substitutes --- the composition C5. Concretely:

\begin{itemize}
\item \textbf{A matched persistence loss with a pair-frozen backward}
  (\S\ref{sec:method}): exact persistence forward (GUDHI alpha complex);
  backward, each paired birth/death simplex re-expressed as the
  closed-form circumradius of live surface samples --- Gabriel-exact for
  100\% of significant pairs on our shapes, so the frozen gradient is
  the true gradient almost everywhere.
  % src: PHASE3_PLAN.md Appendix A (22/22 gate)
\item \textbf{A recruitment term for missing features}: optimal partial
  matching yields exactly zero gradient when a target feature has no live
  counterpart --- the most important failure mode; recruitment attaches
  the nearest live bar instead. A controlled toy probes its known
  location-blindness; removing it in-loop collapses the loop-class win
  to baseline (\S\ref{sec:results}).
\item \textbf{One calibrated knob}: $\lambda$ is set by a gradient-norm
  ratio $\rho$ against the photometric loss; the response is flat over a
  $10\times$ range and a curriculum adds nothing --- an ablation, not an
  assumption.
  % src: PHASE3_PLAN.md Appendix C (rho sweep), Appendix D (C3)
\item \textbf{Channel-controlled evaluation}: every claim is tested against
  a norm-matched \emph{non-topological} control acting through the identical
  channel (a repulsion on the same surface samples) --- the allocation
  study's control discipline, inherited.
\end{itemize}

The loss is topology-specific for enclosed voids ($4.0$--$7.9\times$ at
Chamfer parity) and, unlike every prior we tested, for loops
($2.3\times$, zero phantom handles) --- RH2 confirmed on the class
allocation never won. The channels compose ($7.6\times$; RH3);
components (H0) are again matched by the generic control --- three
channels agree the class carries no exploitable signal, sharpening RH1:
topological guidance earns its keep where geometry and topology
decouple.
% src: PHASE3_PLAN.md Appendix D / topo3/report/results.json
